\section{Methods}
\label{sec:methods}

\subsection{Orbital and Surface Planetary Magnetic Field Datasets}
Planetary magnetic field modeling relied on calibrated satellite and surface magnetometer observations:
\begin{enumerate}
    \item \textbf{Lunar Datasets:} Lunar crustal magnetic field maps at \SI{30}{\kilo\meter} altitude were obtained from the NASA Planetary Data System (PDS) archive \citep{Hood2020, Hood2022}, compiled from vector magnetometer data collected by the Lunar Prospector (LP-MAG) and SELENE (Kaguya) missions. Equatorial and mid-latitude coverage ($\SI{65}{\degree}\text{S}$ to $\SI{65}{\degree}\text{N}$) at $\sim\SI{0.5}{\degree}$ resolution was integrated with high-resolution polar models ($\SI{60}{\degree}$ to $\SI{90}{\degree}\text{S}$) \citep{Hood2022} and surface anomaly models \citep{Tsunakawa2015}.
    \item \textbf{Martian Datasets:} Martian crustal magnetic field determinations were derived from the continuous spherical harmonic model of \citet{Langlais2019} (DOI: 10.1029/2019JE005979), combining vector magnetometer observations from the Mars Global Surveyor (MGS MAG/ER) at \SI{400}{\kilo\meter} mapping orbit and periapsis passes ($\sim\SI{130}{\kilo\meter}$) with low-altitude MAVEN nightside observations. Ground-level surface field magnitudes ($|\mathbf{B}_{\text{surf}}|$) were estimated via downward-continuation scaling and surface magnetic field extrapolation validated against direct in situ surface magnetometry from the NASA InSight fluxgate magnetometer \citep{Johnson2020}.
\end{enumerate}

Total magnetic field intensity ($|\mathbf{B}|$) was calculated from orthogonal spatial components:
\begin{equation}
|\mathbf{B}| = \sqrt{B_r^2 + B_\theta^2 + B_\phi^2}
\end{equation}
where $B_r$ is radial, $B_\theta$ is colatitudinal (southward), and $B_\phi$ is azimuthal (eastward).

\subsection{Landing Site Extraction and Bilinear Interpolation}
A unified database of 36 planetary landing coordinates was constructed, encompassing 23 lunar sites (13 Artemis III candidate South Polar regions \citep{NASA2022}, Apollo 11, 12, 14, 15, 16, 17, Chang'e 3, 4, 5, and Chandrayaan-3) and 13 Martian sites (Arcadia Planitia, Deuteronilus Mensae, Terra Sirenum, Perseverance, Curiosity, InSight, Spirit, Opportunity, Phoenix, Viking 1/2, Zhurong, ExoMars). Local vector magnetic field components were extracted using 2D bilinear spatial interpolation:
\begin{equation}
f(\lambda, \phi) \approx \sum_{i=1}^2 \sum_{j=1}^2 w_{ij} f(\lambda_i, \phi_j)
\end{equation}
where $w_{ij}$ represent inverse-distance spatial weighting factors across latitude $\lambda$ and longitude $\phi$.

\subsection{Systematic Review Protocol}
A systematic review of hypomagnetic field biological effects was conducted across PubMed, Web of Science, and NASA Technical Reports (2000--2026). Search queries combined terms ``hypomagnetic field'', ``near-null magnetic field'', and ``zero magnetic field'' with target biological endpoints (``Arabidopsis'', ``cryptochrome'', ``auxin'', ``ROS'', ``microbiome'', ``osteoclast'', ``bone demineralization'', ``radical pair'', ``DNA repair''). Peer-reviewed experimental studies utilizing Mu-metal or active Helmholtz coil shielding achieving field attenuation below \SI{5}{\micro\tesla} were prioritized for quantitative synthesis.

\subsection{Software and FAIR Data Availability}
Data processing pipelines were executed in Python using \texttt{numpy}, \texttt{pandas}, and \texttt{scipy}. LaTeX tables were generated automatically using custom scripts. The complete codebase, processed datasets, and publication source files are archived on Zenodo (DOI assigned upon publication) and hosted at \url{https://github.com/dr-richard-barker/lunar-mars-magnetic-biology}.
